\section{Designmål og begrensninger}
\label{sec:designmal}

Designet av VideoAPI tok utgangspunkt i kravspesifikasjonen fra HDO. Løsningen
skulle fungere som en mellomvare mellom eksisterende videoløsninger og
videomotorene AntMedia Server og NHN Video. Målet var å tilby ett felles API som
kunne rute kall til valgt backend, uten at klienten måtte kjenne detaljene i hver
videomotor~\cite{hdo_kravspec}. Dette gjorde det nødvendig å redusere koblingen
mellom klientene og de konkrete videomotorene. AntMedia og NHN har ulike API-er
og domenemodeller: AntMedia eksponerer REST-endepunkter for blant annet
broadcast-ressurser, mens NHN Booking API er orientert rundt opprettelse og
administrasjon av videomøter~\cite{antmedia_rest,nhn_booking}. Arkitekturen måtte
derfor skille felles logikk fra backend-spesifikk logikk.

Utvidbarhet var et annet sentralt designmål. Kravspesifikasjonen peker på at nye
videomotorer bør kunne legges til uten at hele mellomvaren må skrives
om~\cite{hdo_kravspec}. Dette ledet til en provider-basert arkitektur, der hver
videomotor implementeres som en egen provider bak et felles grensesnitt. Samtidig
var sikkerhet og driftbarhet sentrale begrensninger: API-et måtte beskyttes med
tokenbasert autentisering, tilgang til videorom måtte kontrolleres med
backend-tilpassede tilgangsmekanismer, og løsningen måtte støtte
containerisering, helsesjekker, logging og Prometheus-metrikker~\cite{hdo_kravspec}.

Et eget designvalg var å holde mellomvaren stateless. Kravspesifikasjonen åpner
for Redis eller database dersom delt state er nødvendig, men i denne løsningen
lagres ikke videorom eller sesjonsdata i VideoAPI. Slik tilstand eies i stedet
av de eksterne videomotorene, mens konfigurasjon hentes fra miljøvariabler i
tråd med krav~I1.2~\cite{hdo_kravspec}. Designet ble dermed styrt av fire
hovedhensyn: et felles API for klientene, isolering av backend-spesifikk logikk,
sikker tilgangskontroll og en driftsvennlig, utvidbar arkitektur.